\problemstart{AP-01-004}{Flattening: Shape Is Not Ownership}{Difficulty 3/5}{Chapter 1 \textbullet\ numpy}{Implement one flattening API whose caller can request aliasing efficiency or independent storage.}

        \subsection{Mathematical statement}


Flattening defines an index bijection
\[
  F(X)_{i,,28r+c}=X_{i,r,c},
  \qquad
  \begin{aligned}
  &0\le r,c<28,\\[-0.2em]
  &F(X)\in\mathbb R^{N\times 784}.
  \end{aligned}
\]
The value equation does not determine the storage relation.  Introduce a policy
$q\in\{\mathsf{shared},\mathsf{independent}\}$ and satisfy both the value map and
the requested ownership relation whenever the input is C-contiguous.


        \subsection{Code problem}


Implement \code{flatten\_images}.  Validate a three-dimensional $28\times28$
image batch.  With \code{require\_independent=False}, a C-contiguous input must
produce a zero-copy view.  With it set to true, the output must never share
memory.  A non-contiguous input must still produce correct row-major values;
\newline Document that NumPy may need a copy in that case.


        \begin{contractbox}
        \begin{Code}
        def flatten_images(images: np.ndarray, *, require_independent: bool) -> np.ndarray:
        \end{Code}
        \end{contractbox}

        \subsection{Theory--engineering gap inventory}

        \begin{gapbox}
        \textbf{Explicit gap.} A bijection of indices specifies values but not strides, contiguity, base objects, mutation propagation, or allocation cost.

        \textbf{Implicit gap exposed by the result.} The exact alias truth table and a reversed-stride fixture show why matching shape and values is insufficient evidence.
        \end{gapbox}

        \subsection{Acceptance evidence}

        \begin{evidencebox}
        \begin{itemize}
          \item Output values match the row-major index equation for contiguous and reversed-stride inputs.
  \item The alias matrix is exactly: contiguous/shared $\mapsto$ true; contiguous/independent $\mapsto$ false.
  \item The benchmark prints storage bytes and mutation propagation rather than assuming \code{reshape} always returns a view. (Trust-based: the test suite does not machine-check benchmark output.)
        \end{itemize}
        \end{evidencebox}

        \subsection{Constraints}

        \begin{itemize}
          \item Use NumPy reshape/copy facilities; do not write index loops.
  \item Do not call \code{flatten()}, which always copies and erases the policy distinction.
  \item Reject with \code{ValueError} a batch that is not three-dimensional or whose trailing dimensions are not $(28,28)$.
        \end{itemize}

        \begin{sourcebox}
        This is an atomic adaptation, not copied solution code. Nearest primary sources:
        \begin{itemize}
          \item \href{https://numpy.org/doc/stable/user/basics.copies.html}{NumPy: Copies and views}
  \item \href{https://d2l.ai/chapter_preliminaries/ndarray.html}{D2L: Data Manipulation}
        \end{itemize}
        \end{sourcebox}
